\section{The original finite holomorphic CM uniformization family}

\paragraph{Scope and conventions.}
This chapter proves the missing family theorem, rather than assuming it
from its individual-point consequence.  Throughout, \(E,H,D,l,X_o\)
are the original data, with \(X_o:\texttt{Moduli}\ D\ l\) supplied.
Put \(R=\mathcal O_E\), \(f=[E:\mathbb Q]/2\), \(r=p+q\geq1\),
\(g=fr\), \(d=fpq\), and \(N=l.M'=\operatorname{lcm}(M,A)\).
The fixed map \(\sigma:K_A(l)\to\mathbb C\) satisfies
\(\sigma(\rho)=\mathrm{rt}_0\).  The \emph{given ordered} type
\(\mathrm{ty}:\operatorname{Fin}(g)\to\operatorname{Hom}(E,\mathbb C)\)
has multiplicities \(p\) on \(\Psi\), and \(q\) on \(\bar\Psi\).
We retain every original level condition, including odd \(M\geq3\),
the two relatively prime half-ideals, and arbitrary \(A>0\).
Hermitian forms are complex linear in their first variable.

Write \(L=D.\mathrm{La}\), \(B=D.\mathrm{EF}\).
For another datum \(D'\) write \(L'\) and \(B'\).
The notation
\[
 \mathcal U(D',t,P;y,e)
 \tag{HUF.1}
\]
means the \emph{literal} \texttt{Unif} predicate in
\texttt{FamilyUnifOfParts}: \(y\) is over \(\sigma\);
\(P:L'\to\mathbb C^g\) is an original \texttt{IsPeriod};
\(e:\mathbb C^g\to G_y(\mathbb C)\) is a surjective holomorphic
group map with kernel \(P(L')\), with both source dual-number
derivative bijections and the exact diagonal \(R\)-action; and
\[
 \begin{split}
 e_n(e(Pa/n),e(Pb/n))&=\exp(-2\pi i B'(a,b)/n)\quad(n>0),\\
 \eta(a)&=e(P(ta)/M)\quad(a\in\bar{\mathfrak b}L),\\
 \eta_A(a)&=e(P(ta)/A)\quad(a\in L).
 \end{split}                                                     \tag{HUF.2}
\]
The first line is the source's universally tested
\texttt{Out.WeilIs}, not an unspecified pairing convention.
All functions may be assigned arbitrary values outside the open sets
on which their conditions are imposed.

\paragraph{Inherited results and named foundations.}
We use PUN6, PUN8 and PUN15, and the actual individual-point,
line-normalization, equality and nonemptiness proofs PUN20--PUN24.
In particular the original primitive identity fixes the
Appell--Humbert semicharacter to
\(\alpha(Pa)=\exp(\pi i B'(a,w_0a))\).
The supplied \texttt{Moduli} record retains its algebraic representing
property on reduced schemes, smoothness of relative dimension \(d\),
quasiprojectivity, universal object and stated infinitesimal condition.
Existence of that record is not proved here.

The retained classical foundations are Appell--Humbert and its
positivity criterion, the Riemann algebraicity theorem, projective
GAGA, elementary holomorphic differential equations and the
implicit-function theorem.  Two additional precise references are
used below: Grauert's analytic direct-image and base-change theorem
for a proper map, a sheaf flat over a reduced base, and constant
fiber cohomology dimension; and finite connected-component
decomposition for semialgebraic sets.  These are external foundations,
not numbered claims of a newly written proof of those classical
theorems.  Exact primary-source locators appear after their applications
and in the component metadata.

\subsection{Reverse extraction on the original moduli space}

% BEGIN HUF1
\paragraph{HUF1. The source holomorphic and tangent tests are the analytic ones.}
\label{claim:HUF1}
Let \(Z\) be a finite-type scheme over a subfield \(k\subseteq\mathbb C\),
with a fixed embedding \(k\to\mathbb C\), and let
\(Z_\mathbb C\) denote the corresponding base change.
For a map \(y:B\to Z_\mathbb C(\mathbb C)\), where
\(B\subseteq\mathbb C^m\) is open, the original
\(\texttt{IsHolTo}(B,y)\) is equivalent to holomorphy into
\(Z_\mathbb C^{\mathrm{an}}\).
At a smooth point, the original \texttt{IsDerivAt0} is the
usual derivative interpreted as a map from the dual numbers.
Consequently a holomorphic local biholomorphism at \(0\) has
both of the exact unique-existence clauses in \texttt{ex\_diff}.

\begin{proof}
On an affine open choose finitely many algebra generators and its
closed immersion into affine space.  The source test says that the
inverse image of this open is open and that all those coordinate
functions are holomorphic.  Their polynomial equations hold pointwise,
so this is precisely a holomorphic map into the closed analytic
subspace.  Coefficients introduced by base change are complex
constants; they cause no additional test.  Conversely a holomorphic
map has holomorphic pullbacks of regular functions.  This remains true
on any scheme open: locally its regular functions are quotients of
polynomials with nonzero denominator.  These local statements glue
and give the quantifier over every open and every section in
\texttt{IsHolTo}.

A morphism from \(\operatorname{Spec}\mathbb C[\epsilon]/(\epsilon^2)\)
with specified closed point is, on its local ring, exactly
\(a\mapsto a(y)+\epsilon\,\delta(a)\), with \(\delta\) a complex
derivation.  Holomorphic directional differentiation gives that
derivation.  Smooth algebraic coordinates, obtained from a standard
smooth presentation and the analytic implicit-function theorem, make
this an isomorphism between analytic and algebraic tangent spaces.
Evaluating every local regular function is exactly the equality
required in \texttt{IsDerivAt0}; restricting an open does not change
the germ or its derivative.  If the analytic derivative is invertible,
each direction determines a unique dual-number point, and each
vertical dual-number point determines a unique direction.
For a group scheme over \(y\) the condition of being vertical is
exactly the source's \texttt{epsPt} condition.
\end{proof}
% END HUF1

% BEGIN HUF2
\paragraph{HUF2. The actual fibers and their relative Lie eigenspaces.}
\label{claim:HUF2}
The complex base change of the supplied universal group is a proper
smooth family of connected abelian varieties of dimension \(g\).
Its Lie bundle is the direct sum of the original embedding
eigenbundles, of ranks \(p\) and \(q\) in the order prescribed
by \(\mathrm{ty}\).  Locally on an analytic base, those eigenbundles
have simultaneous holomorphic frames with any prescribed eigenbasis
at one point.  The universal rigidified line is relatively ample
and its analytic restriction to each fiber has a positive integral
Appell--Humbert form.

\begin{proof}
The \texttt{SAB.semi} field gives a smooth group with connected
semiabelian geometric fibers.  The additional properness field
makes its toric part trivial: a positive-dimensional torus cannot
be a closed subgroup of a proper group, since it is affine and
not proper.  Thus the fibers are abelian varieties.  The
\texttt{Obj.geom.DetCond} clause computes their Lie multiplicities.
Summing them over the embeddings gives \(g=f(p+q)\);
smoothness with this fiber dimension gives the asserted relative
dimension.

Choose an integral primitive element for \(E/\mathbb Q\).
Its distinct conjugates allow the usual interpolation polynomials
in its Lie action to construct the embedding idempotents.
They are holomorphic bundle endomorphisms.  Their images are
holomorphic subbundles of the constant ranks given by
\texttt{DetCond}.  Choose local frames and multiply them by constant
invertible matrices to match the prescribed basis at the initial
point.  Ordering these frames according to the supplied type gives
the diagonal \(R\)-action, without permuting the original coordinates.

The printed \texttt{RelAmple} fields have a useful concrete meaning.
On every affine base open they give a power of the line, finitely
many generating sections, and affine nonzero loci whose coordinate
rings are generated by the section ratios.  The associated map
to projective space is therefore a closed immersion: on each
standard projective open its coordinate-ring map is surjective.
In particular its analytification is proper, since it is closed
in the product with the compact projective space.  The same
construction after base change gives the proper analytic family
used here.  The line's rigidification and primitive, symmetric
and adjoint identities are identities of the actual universal
object and hence persist.  On a fiber, Appell--Humbert and the
positivity criterion identify the ample line with a positive
Hermitian form and an integral alternating imaginary part.
\end{proof}
% END HUF2

% BEGIN HUF3
\paragraph{HUF3. Exponential, full period lattice and the original Weil sign.}
\label{claim:HUF3}
Let \(G\) be any of the complex fibers in HUF2.
There is a holomorphic exponential
\[
       \exp_G:\operatorname{Lie}G\longrightarrow G(\mathbb C)
 \tag{HUF.3}
\]
with derivative the identity, surjective with discrete full
lattice kernel \(\Lambda\).  For an Appell--Humbert form
\(H_G\) of its line, put \(B_G=\operatorname{Im}H_G|_\Lambda\).
Then \(B_G\) is a perfect integral alternating form.  In the
original commutator convention its Weil scalar is
\[
 e_n(\exp_G(\lambda/n),\exp_G(\mu/n))
                     =\exp(-2\pi iB_G(\lambda,\mu)/n).
 \tag{HUF.4}
\]
The conclusion includes the universal condition on generators
in \texttt{WeilIs}.

\begin{proof}
Left-invariant holomorphic vector fields commute because the group
is commutative.  Their local flows give a local holomorphic
homomorphism with identity derivative.  For an arbitrary vector
\(z\), choose a positive integer \(a\) with \(z/a\) in its
domain and define \(\exp_G(z)=[a]\exp_G(z/a)\).
The local homomorphism law makes the choices compatible, by
passing to a common multiple.  It also proves global
holomorphy and addition.  Its image is an open subgroup,
and connectedness forces it to be the whole group.
The derivative is invertible everywhere by translation, so the
kernel is discrete.  The induced local biholomorphism
\(\operatorname{Lie}G/\Lambda\to G(\mathbb C)\) is bijective
and hence an analytic isomorphism.  Compactness of \(G(\mathbb C)\)
makes \(\Lambda\) a full lattice: a discrete subgroup of a real
vector space has a basis for its real span, and a proper span
would leave a noncompact real-vector-space quotient.

Appell--Humbert supplies the integral form \(B_G\).
Its positivity makes it nonsingular.  The kernel of the line's
polarization consists analytically of the quotient of the
integral \(B_G\)-dual lattice by \(\Lambda\).  This follows
directly by comparing the factors for a translate of the line:
the translation difference is trivial exactly for a vector
pairing integrally with \(\Lambda\).  Such a trivial analytic
line is algebraically trivial by projective GAGA.
The source \texttt{KTriv}, applied with coefficient ring
\(\mathbb C\), forces the corresponding point to be zero.
Thus the dual lattice is \(\Lambda\), proving perfection.

For \(x=\lambda/n\), the section of the source's translation
quotient is represented by
\(a_x(z)=\exp(-\pi nH_G(z,x))\).
The semicharacter cancels from that quotient.  Its factors
are exactly those required to descend this nonvanishing
section.  As in the explicit calculation PUN21, the source
commutator, in its printed order, is
\[
 \frac{a_x(z)a_y(z+x)}{a_y(z)a_x(z+y)}
   =\exp[-\pi n(H_G(x,y)-H_G(y,x))]
   =\exp(-2\pi iB_G(\lambda,\mu)/n).
\]
GAGA makes these sections algebraic.  Every other generator
differs by a nonzero constant on the connected proper fiber,
so the same scalar holds for every generator.  This proves
the literal \texttt{WeilIs}, not only the value for the
exhibited analytic representatives.  HUF1 supplies the two
derivative clauses for (HUF.3).
\end{proof}
% END HUF3

% BEGIN HUF4
\paragraph{HUF4. Recovering the original integral CM datum from a fiber.}
\label{claim:HUF4}
For every \(y\) over the fixed \(\sigma\), the exponential
kernel of \(G_y\) carries an original datum \(D'\) of
signature \((p,q)\), and an original ordered period
\(P:L'\to\mathbb C^g\).  Its exponential has the exact
diagonal \(R\)-action, form \(B'\), tangent dictionary and
all Weil pairings in (HUF.1)--(HUF.2).

\begin{proof}
Differentiate each \(R\)-endomorphism.  Uniqueness of local
flows, followed by division by integers, gives
\[
          \exp_G(d[a]\,z)=[a]\exp_G(z).
 \tag{HUF.5}
\]
Thus the derivative action preserves the actual lattice
\(\Lambda\).  The adjoint identity of the actual line, read
in its mixed Appell--Humbert Hermitian forms, gives
\(B_G(a\lambda,\mu)=B_G(\lambda,\bar a\mu)\).
The lattice is finite over \(R\), since its finite
\(\mathbb Z\)-basis already generates it over \(R\).
It is torsion free: a nonzero \(a\in R\) acts on each Lie
eigenspace by the nonzero number \(\tau(a)\).
Over the Dedekind domain \(R\), a finite torsion-free module
is projective.  Its rank is \(r\), because its
\(\mathbb Z\)-rank is \(2g=[E:\mathbb Q]r\).

Use the given ordered type to identify the Lie algebra
complex linearly with \(\mathbb C^g\), and let \(P\) be
the image of the inclusion of \(\Lambda\).
Transport multiplication by \(i\) to \(J\) on
\(\mathbb R\otimes\Lambda\).  The form \(B_G\) is
\(J\)-invariant and \(B_G(Jv,v)=H_G(v,v)>0\) for
\(v\ne0\).  The exact \texttt{Datum.dom} trace identities
follow on each embedding eigenspace:
\[
 \operatorname{tr}_{\mathbb R}a
   =2\operatorname{Re}\sum_{\tau\in\Psi}
                       (p\tau(a)+q\bar\tau(a)),\qquad
 \operatorname{tr}_{\mathbb R}(Ja)
   =-2\operatorname{Im}\sum_{\tau\in\Psi}
                       (p\tau(a)+q\bar\tau(a)).
\]
Perfection is HUF3.  These verifications give every field
of \(D'\), including its original integral lattice, rather
than a lattice of finite index.  Discreteness, full real
span, positivity and the diagonal action prove
\texttt{IsPeriod}.  Equations (HUF.4)--(HUF.5) and HUF1
give the remaining stated exponential fields.
\end{proof}
% END HUF4

% BEGIN HUF5
\paragraph{HUF5. Exact comparison with the original geometric Tate marking.}
\label{claim:HUF5}
The construction of HUF4 admits an \(R\)-linear map
\(t:L\to L'\) for which
\(\texttt{InGenus}(D,l,\sigma,D',t)\) holds and the two
marking identities in (HUF.2) hold.  In particular all
the compatible maps \(f_n\) and multipliers \(\nu_n\)
in the original \texttt{InGenus} are produced.

\begin{proof}
Use the \texttt{Obj.geom} clause at \(y\), with coefficient
field \(\mathbb C\) and the fixed \(\sigma\).  It provides
the original \(\texttt{TateIso}\ I\), its chosen roots of
unity, and its two marking identities.
The exponential provides the second Tate description
\[
  I'_n(a)=\exp_G(Pa/n),\qquad
                      \zeta'_n=\exp(-2\pi i/n)\quad(n>0).
\]
Its kernel is \(nL'\); every \(n\)-torsion point has this
form, by surjectivity and the exact exponential kernel.
Its tower law and \(R\)-linearity follow from (HUF.5),
and its pairing is HUF3.

For each \(n>0\), comparison of the two descriptions gives
an \(R\)-linear isomorphism
\(\bar f_n:L/nL\to L'/nL'\).
The composed map \(L\to L'/nL'\) lifts to an
\(R\)-linear \(f_n:L\to L'\), because the \emph{given}
\(L\) is projective over \(R\).  This is an integral
lift through the quotient \(L'\to L'/nL'\); no choice of
an \(R\)-basis is required.  The torsion tower laws give
\(f_{nm}(a)-f_n(a)\in nL'\) for all \(n,m>0\).
Each \(\bar f_n\) is an isomorphism, giving exactly the
two source injectivity and surjectivity clauses modulo \(n\).

Choose an integer \(\nu_n\) such that
\(I.\zeta_n=\exp(-2\pi i\nu_n/n)\).  Primitivity makes
\(\nu_n\) invertible modulo \(n\).
Comparing the two Weil values, whose values are unique,
gives
\[
           B'(f_na,f_nb)-\nu_nB(a,b)\in n\mathbb Z.
\]
The specified root at \(A\) gives the required
\(\sigma(\zeta_A)=\exp(-2\pi i\nu_A/A)\).
Set \(t=f_N\).  Since \(M,A\mid N\), compatibility
gives \(f_N\equiv f_M\bmod M\) and
\(f_N\equiv f_A\bmod A\).  Substituting in the geometric
Tate marking identities gives exactly (HUF.2), on the
source's half-lattice for \(M\) and on all of \(L\)
for \(A\).  Values \(f_0=0,\nu_0=0\) are harmless:
all the relevant genus congruences have their original
positive-\(n,m\) guards.
\end{proof}
% END HUF5

% BEGIN HUF6
\paragraph{HUF6. The whole reverse-uniformization-at-a-point producer.}
\label{claim:HUF6}
For every original \(E,H,D,l,X_o,\sigma,g,\mathrm{ty}\)
in the source scope and every \(y\) over \(\sigma\),
there exist \(D',t,P,e\) such that
\[
 \texttt{InGenus}(D,l,\sigma,D',t)
                 \quad\hbox{and}\quad \mathcal U(D',t,P;y,e).
 \tag{HUF.6}
\]
This is the literal
\texttt{PointUnif.exists\_uniformisation\_at\_point\_closed}
conclusion, with the original moduli point unchanged.

\begin{proof}
Pull the supplied universal object back to \(y\).
HUF2 identifies its actual abelian variety and line.
HUF3 constructs the exponential with the required
holomorphic and derivative fields.
HUF4 extracts the integral datum and the period with the
given ordering of the type, and HUF5 supplies the exact
genus and both markings.  These are all conjuncts of
(HUF.6).  Neither existence of the moduli scheme nor
existence of a family uniformization has been used.
The fixed embedding \(\sigma\), including its chosen
root of unity, is the one used in HUF5.
\end{proof}
% END HUF6

\subsection{Prescribed local lifts and holomorphic dependence}

% BEGIN HUF7
\paragraph{HUF7. Relative exponentials and continuation of the full lattice.}
\label{claim:HUF7}
For a holomorphic family \(y:B\to X_{o,\sigma}^{\mathrm{an}}\)
and \(b_0\in B\), fix an exponential at \(b_0\) with its
given ordered Lie identification and an integral basis
of its full period lattice.  On a sufficiently small
connected open polydisc about \(b_0\) there is a jointly
holomorphic relative exponential, and the prescribed
period basis extends holomorphically to a basis of its
\emph{entire} kernel at every point.

\begin{proof}
Pull back the proper analytic group of HUF2 to \(B\).
Choose the local eigenframes of HUF2 to agree with the
prescribed initial derivative.  In these frames the
relative invariant vector fields depend holomorphically
on both their vector and base parameters.  The local
holomorphic differential-equation theorem gives a local
exponential near the zero section.  Define it at arbitrary
vectors by multiplication after division by a sufficiently
large positive integer, exactly as in HUF3.  On any
neighborhood of a given \((b,z)\) the same integer works.
The compatible formulas therefore give a jointly
holomorphic relative exponential on \(B'\times\mathbb C^g\).

Its derivative in the vector direction is invertible
everywhere.  The analytic implicit-function theorem,
applied to the equation \(\exp_b(\lambda)=0\) near each
of the finitely many prescribed basis vectors, continues
them as holomorphic kernel sections \(\lambda_j(b)\).
Shrink once so that their real coefficient matrix remains
invertible.  They generate a lattice \(\Lambda_b^0\).
The action of \(\mathbb Z^{2g}\) by these translations
on \(B'\times\mathbb C^g\) is properly discontinuous:
on every compact base subset the inverse real coefficient
matrices have a common bound.  The quotient
\[
       T^0=(B'\times\mathbb C^g)/\Lambda^0
                         \longrightarrow B'
\]
is a holomorphic torus family, topologically the product
with \(\mathbb R^{2g}/\mathbb Z^{2g}\), and is proper.
The exponential descends to a map \(T^0\to G_y\).
It is a local biholomorphism, since its vertical derivative
is invertible and it is the identity on the base.
On every fiber it is onto: its image is an open subgroup
of a connected compact group.

This map is proper.  Indeed the inverse image of a
compact subset is closed in the inverse image in \(T^0\)
of its compact projection to \(B'\).  A proper local
biholomorphism is a finite covering.  Its degree is
constant on the connected total space \(G_y|_{B'}\).
At \(b_0\) that degree is one, because the initial
period basis was the full kernel.  Thus it is an
isomorphism, and \(\Lambda_b^0\) is the full kernel
for every \(b\).  This proves fullness; continuation
of a finite collection of kernel elements alone would
not have sufficed.
\end{proof}
% END HUF7

% BEGIN HUF8
\paragraph{HUF8. One neighborhood preserves the form, action and both markings.}
\label{claim:HUF8}
Suppose the initial data in HUF7 satisfy
\(\mathcal U(D',t,P_0;y(b_0),e_0)\).
The continued lattice gives, on one open \(B'\ni b_0\),
\[
 P(b_0)=P_0,\qquad \texttt{IsHolPeriods}(D',\mathrm{ty},B',P),
 \qquad
 \mathcal U(D',t,P(b);y(b),e_b)\quad(b\in B').
 \tag{HUF.7}
\]
The maps \(e_b\) are jointly holomorphic.  This is the
whole original \texttt{exists\_local\_period\_lift}, including
preservation of its prescribed initial period.

\begin{proof}
Choose a \(\mathbb Z\)-basis of \(L'\) and continue its
images under \(P_0\) as in HUF7.  Extend by integral
linearity to obtain \(P(b)\).  The chosen eigenframes
ensure the desired diagonal \(R\)-action on the Lie
bundle.  Applying an element of \(R\) to a continued
period gives another continuous kernel section.  Its
integer coordinates in the continued basis are locally
constant, so are their original coordinates at \(b_0\).
This proves \(R\)-linearity for a finite integral basis
of \(R\), and hence for every element.  The period map
at \(b_0\) is the original \(P_0\), and uniqueness of
the exponential with a specified derivative gives
\(e_{b_0}=e_0\).

Pull the actual line back to the torus family of HUF7.
Its integral Appell--Humbert alternating form is locally
constant, because it is determined, with our fixed sign
convention, by the first Chern class on the continued
two-cycles.  One can see this
without choosing a connection: identify the underlying
torus topologically with a fixed torus on the polydisc;
the restrictions of a line bundle along two homotopic
fiber inclusions have the same integral Chern class.
The resulting alternating matrix is therefore \(B'\)
at every point.  At the initial point this follows
also from the full Weil identities: the difference
of its actual matrix and \(B'\) is divisible by
every positive integer.  Positivity of the ample line
supplies the positive Hermitian matrix at each fiber.
Thus every \(P(b)\) is precisely an \texttt{IsPeriod}.

The two markings are torsion sections, of orders dividing
\(M\) and \(A\).  Multiplication by either positive
integer on the analytic family is a local biholomorphism
with finite fibers, so its kernel is a finite covering.
Two continuous sections of this kernel that agree at
one point agree throughout a connected small base.
Apply this to a finite integral generating set of
\(\bar{\mathfrak b}L\) for the first marking, and of
\(L\) for the second.  Addition gives both equalities
for every required vector.  Finally HUF3 computes
all positive-\(n\) Weil pairings from the single
constant form \(B'\); it does not require one
neighborhood for each \(n\).  HUF1 gives the tangent
clauses.  These arguments establish all of (HUF.7)
on one neighborhood.  This proof needs no
\texttt{InGenus} hypothesis beyond the stated initial
\(\mathcal U\), exactly as in the source local-lift
theorem.
\end{proof}
% END HUF8

% BEGIN HUF9
\paragraph{HUF9. Local period coordinates have exactly dimension \(fpq\).}
\label{claim:HUF9}
For the original \(D'\), every ordered period \(P_0\)
has a finite-period neighborhood and a holomorphic
coordinate \(Z(P)\in\mathbb C^{fpq}\) with the following
property: if \(P,P'\) are in that neighborhood, then
\[
 Z(P)=Z(P')\ \Longrightarrow\
 \exists h\in\operatorname{GL}_g(\mathbb C):
          P'=hP,\quad h\text{ commutes with the diagonal }R\text{-action}.
 \tag{HUF.8}
\]
For every holomorphic period family in the neighborhood,
\(Z(P(b))\) is holomorphic.  This proves the exact
\texttt{PeriodCoord.exists\_period\_coordinates} interface.
As a standalone assertion it also includes \(p+q=0\):
then \(g=0\), the lattice is zero and the coordinate
space is a point.

\begin{proof}
Put \(W=L'\otimes_{\mathbb Z}\mathbb C\) and extend \(P\)
complex linearly to \(W\to\mathbb C^g\).
Its kernel \(Q(P)\) is a complex \(g\)-plane.
The full real span implies
\(Q(P)\oplus\overline{Q(P)}=W\).
The Riemann identities say that \(B'_\mathbb C\) vanishes
on \(Q(P)\), and positivity is an open condition.
Decompose \(W=\bigoplus_\tau W_\tau\).
Each \(W_\tau\) has dimension \(p+q\).
The adjoint relation makes \(B'_\mathbb C\) pair
\(W_\tau\) perfectly with \(W_{\bar\tau}\) and with
no other summand.

For \(\tau\in\Psi\), the kernel \(Q_\tau\) has dimension
\(q\).  Its partner \(Q_{\bar\tau}\), of dimension \(p\),
is its annihilator under this perfect pairing.
Thus \(Q\) is specified by the \(q\)-planes \(Q_\tau\)
for \(\tau\in\Psi\), subject only to the open
real-complementarity and positivity conditions.
It lies in an open subset of
\(\prod_{\tau\in\Psi}\operatorname{Gr}_q(\mathbb C^{p+q})\),
of complex dimension \(fpq\).
Choose complements at \(Q(P_0)\).  Nearby kernels are
graphs of \(q\)-by-\(p\) matrices, which give \(Z\).
Their entries are quotients of matrix minors of \(P\)
with fixed nonzero denominators at \(P_0\); hence they
are holomorphic functions of finitely many period
coordinates.

Choose a finite integral basis of \(L'\).
Continuity of those finitely many minors gives an
\(\varepsilon>0\) such that
\(\|P(a)-P_0(a)\|<\varepsilon\) on that basis places
\(P\) in the graph chart.  This supplies the finite
set and the common radius in the literal source
statement, for every parameter dimension \(m\).
Equal graph coordinates mean equal kernels.  The
two quotient isomorphisms \(W/Q\to\mathbb C^g\)
then differ by a unique \(h\), proving (HUF.8).
Both maps are \(R\)-linear, so \(h\) commutes with
the prescribed action.  If \(p=0\) or \(q=0\),
these Grassmannians are points and the same argument
gives the required zero-dimensional coordinate.
\end{proof}
% END HUF9

% BEGIN HUF10
\paragraph{HUF10. A marked exponential is unique.}
\label{claim:HUF10}
Suppose \(\texttt{InGenus}(D,l,\sigma,D',t)\) and
both \(\mathcal U(D',t,P;y,e)\) and
\(\mathcal U(D',t,P;y,e')\) hold.  Then \(e=e'\).
This is the literal \texttt{exp\_eq\_of\_periods\_of\_marking}.

\begin{proof}
By HUF1 and the derivative clauses, both maps induce
analytic isomorphisms \(\mathbb C^g/P(L')\to G_y\).
Their quotient gives an analytic group automorphism
of this torus.  Lifting to its simply connected
covering and differentiating gives a complex-linear
automorphism \(h\) and an integral lattice
automorphism \(\gamma\), with \(P\gamma=hP\).
The \(R\)-action implies that \(\gamma\) is \(R\)-linear.
Comparison of all the Weil identities gives
\(B'(\gamma a,\gamma b)-B'(a,b)\in n\mathbb Z\)
for every \(n>0\).  Hence \(\gamma\) is an exact
isometry.

The level identities give
\(\gamma(ta)-ta\in ML'\) for
\(a\in\bar{\mathfrak b}L\), and
\(\gamma(ta)-ta\in AL'\) for all \(a\in L\).
The genus map \(t\) is onto modulo \(N\).
Writing an arbitrary element of
\(\bar{\mathfrak b}L'\) as a sum of ideal multiples
and using this surjectivity shows that it is congruent
modulo \(N\bar{\mathfrak b}L'\) to \(ta\) for
some \(a\in\bar{\mathfrak b}L\).  The half-level
congruence therefore holds on that whole half-lattice.
PUN6 turns it into \(\gamma\equiv1\pmod M\)
on all of \(L'\).

The positive Hermitian form is determined by its
imaginary part and the complex structure.  Thus
\(\gamma\) preserves the positive real norm as well.
An integral isometry group of a positive norm is
finite: images of a fixed integral basis lie in
finite intersections of a lattice with bounded
sets.  Consequently \(\gamma\) has finite order.
PUN8's principal-congruence argument, for the
original odd \(M\geq3\), forces \(\gamma=1\).
Full real span then gives \(h=1\), and \(e=e'\).
No prime \(3\) exception is used.
\end{proof}
% END HUF10

% BEGIN HUF11
\paragraph{HUF11. The period coordinate on a moduli chart is a local biholomorphism.}
\label{claim:HUF11}
At any \(y_0\) with a chosen \(\mathcal U(D',t,P_0;y_0,e_0)\),
there is an analytic chart of the original moduli space
on which the coordinate \(Z\) of its continued periods
is a local biholomorphism to an open subset of
\(\mathbb C^d\).

\begin{proof}
The supplied smoothness field and HUF1 give a
holomorphic injective chart \(x:U\subseteq\mathbb C^d
\to X_{o,\sigma}^{\mathrm{an}}\), with \(x(0)=y_0\).
Apply HUF8 to it, with the prescribed \(P_0,e_0\),
and shrink into the finite-period neighborhood of
HUF9.  The map \(z\mapsto Z(P_z)\) is holomorphic.
If two of its values coincide, HUF9 gives a
complex-linear \(h\) taking one period to the other.
PUN23, with integral isometry \(\gamma=1\) and
the two identical marking maps \(t\), gives equality
of the two moduli points.  The chart is injective,
so the two parameters are equal.

An injective holomorphic map between open subsets
of \(\mathbb C^d\) has invertible derivative at
every point and is a biholomorphism onto its image.
We use exactly this classical theorem, as stated
and proved in Lebl, \emph{Tasty Bits of Several
Complex Variables}, Theorem 1.6.6, pp.\ 43--45.
Its equal-dimension hypothesis holds because
HUF9 gives \(fpq=D.\mathrm{dim}\), the dimension
of the supplied smooth moduli chart.
The holomorphic inverse-function theorem gives
the asserted inverse on an open image.
For \(d=0\), both charts are a single point
and the statement is immediate.
\end{proof}
% END HUF11

% BEGIN HUF12
\paragraph{HUF12. Arbitrarily chosen pointwise uniformizations vary holomorphically.}
\label{claim:HUF12}
Retain \(\texttt{InGenus}(D,l,\sigma,D',t)\).
Let \(B\subseteq\mathbb C^m\) be any open set and
\(P(b)\) any original \texttt{IsHolPeriods} family.
Suppose arbitrary functions \(y(b),e_b\) satisfy
\(\mathcal U(D',t,P(b);y(b),e_b)\) for \(b\in B\).
Then
\[
 \texttt{IsHolTo}(B,y),\qquad
 \texttt{IsHolTo}(B\times\mathbb C^g,(b,z)\mapsto e_b(z)).
 \tag{HUF.9}
\]
These are the full original
\texttt{isHolTo\_modulipoint\_exp\_of\_holomorphic\_periods}
conclusions; continuity of the chosen functions is not
an added hypothesis.

\begin{proof}
Fix \(b_0\).  Use \(y(b_0),P(b_0),e_{b_0}\) in HUF11.
For \(b\) close enough to \(b_0\), the holomorphic
coordinate \(Z(P(b))\) lies in the image of that chart.
Its holomorphic inverse produces a moduli point
\(\widetilde y(b)\) and the chart's jointly
holomorphic periods \(\widetilde P(b)\) and exponential
\(\widetilde e_b\).  Equal \(Z\)-coordinates give
an \(R\)-equivariant complex-linear isomorphism
\(h(b)\) with
\(\widetilde P(b)=h(b)P(b)\).
This matrix is holomorphic: choose \(g\) period
vectors independent over \(\mathbb C\) at \(b_0\);
locally it is the matrix of their tilded columns
times the inverse matrix of their untilded columns.

PUN23, with \(\gamma=1\), gives
\(\widetilde y(b)=y(b)\) for each \(b\), without
first requiring that \(y\) be continuous.
Moreover \(\widetilde e_b\circ h(b)\) and \(e_b\)
have the same complete \(\mathcal U\) data:
the period, action, kernel, tangent, Weil and
marking formulas transport directly under \(h(b)\).
HUF10 therefore identifies them.  The former is
jointly holomorphic in \((b,z)\), so the latter is
as well.  HUF1 gives the precise source predicates.
Doing this at every \(b_0\) proves (HUF.9) on
the whole possibly disconnected open \(B\).
\end{proof}
% END HUF12

\subsection{Finite disjoint families on the supplied moduli scheme}

% BEGIN HUF13
\paragraph{HUF13. Finiteness of the analytic component set.}
\label{claim:HUF13}
The analytic space \(X_{o,\sigma}^{\mathrm{an}}\) has
finitely many connected components.  In particular every
locally constant map from it to a discrete set has finite
image.

\begin{proof}
The supplied quasiprojectivity gives finite type and
quasicompactness, so choose a finite affine cover after
base change to \(\mathbb C\).  Each affine has a closed
immersion into some \(\mathbb C^a\).  In real and
imaginary coordinates its complex points are a real
semialgebraic set: a complex polynomial equation is
two real polynomial equations, with arbitrary real
coefficients allowed.  The finite connected-component
theorem for semialgebraic sets applies.
The finitely many connected pieces of the finitely
many affines cover \(X_{o,\sigma}^{\mathrm{an}}\).
Every such piece lies in one connected component,
so the number of components of the whole space is
at most the number of pieces.  On each connected
component a locally constant map to a discrete set
is constant, proving the assertion.

The named input is Coste, \emph{An Introduction to
Semialgebraic Geometry} (October 2002), Theorem 2.23,
p.\ 40, following the finite cylindrical decomposition
Theorem 2.20.  We use only finiteness of connected
components.  We do not assume that the moduli space
itself is compact, or that its period map is proper.
\end{proof}
% END HUF13

% BEGIN HUF14
\paragraph{HUF14. The exact marked-lattice classes have disjoint open images.}
\label{claim:HUF14}
Consider all pairs \((D',t)\) in the original genus.
Declare \((D',t)\sim(D'',t')\) if there is an
\(R\)-linear isometry \(\gamma:L'\simeq L''\) with
\[
 \begin{split}
 \gamma(ta)-t'a&\in ML''\quad(a\in\bar{\mathfrak b}L),\\
 \gamma(ta)-t'a&\in AL''\quad(a\in L).
 \end{split}                                                     \tag{HUF.10}
\]
The images in \(X_{o,\sigma}(\mathbb C)\) of all periods
for two such pairs are equal if the pairs are equivalent,
and disjoint otherwise.  Each image is nonempty and open,
and together they cover all of \(X_{o,\sigma}(\mathbb C)\).

\begin{proof}
Composition and inversion of \(R\)-linear integral
isometries preserve both congruences in (HUF.10);
thus this is an equivalence relation.
Every datum has periods of the given ordered type
by PUN15.  PUN22 constructs a point for each such
period and genus marking, so each image is nonempty.

If two images meet, apply PUN23 to the two
uniformizations of that same point.  Its integral
isometry has exactly the congruences (HUF.10).
Conversely, given \(\gamma\) as there, a period
\(P\) for \(D'\) transports to \(P\gamma^{-1}\)
for \(D''\).  It has the same exact pairing,
diagonal action, discreteness, fullness and positivity.
PUN22 constructs a point for it, and PUN23 with
complex map \(h=1\) identifies this point with the
one for \(P\).  Thus equivalent pairs have identical
images, not merely an intersection.

At a point in any image, take a small analytic
moduli chart.  HUF8, starting from that point's
uniformization, gives periods on the \emph{same}
\(D',t\) for every nearby point.  The image is
therefore open.  HUF6 proves that all the images
cover the moduli space.  They form a disjoint open
partition, hence define a locally constant class
map to the discrete set of classes.
There is no set-theoretic choice over a proper
class hidden here: one may transport each free
\(\mathbb Z\)-module \(L'\) to
\(\mathbb Z^{2g}\), and then take the set of
its integral actions, forms and marking matrices.
Transporting back gives the original arbitrary
carriers and loses no pair.
\end{proof}
% END HUF14

% BEGIN HUF15
\paragraph{HUF15. One finite positive family, chosen before every period.}
\label{claim:HUF15}
There are \(n>0\) original data \(D_i\) and genus
maps \(t_i:L\to L_i\), one for each of the classes
in HUF14, and total functions
\[
 u_i:(L_i\to^+\mathbb C^g)\to X_o(\mathbb C),\qquad
 e_{i,P}:\mathbb C^g\to G_{u_i(P)}(\mathbb C)
 \tag{HUF.11}
\]
such that \(\mathcal U(D_i,t_i,P;u_i(P),e_{i,P})\)
for every valid period.  These same functions have
the original \texttt{onto}, \texttt{disj},
\texttt{hol} and \texttt{ex\_diff} fields.

\begin{proof}
HUF13 and the locally constant class map of HUF14
make the class image finite.  PUN24 supplies an
actual point over the prescribed \(\sigma\);
HUF6 places it in that image, so its cardinality
\(n\) is positive.
Choose one \((D_i,t_i)\) for each distinct class.
The genus conditions give their correct signatures
and both injectivity and surjectivity of \(t_i\)
modulo \(N\), as required in \texttt{t\_inj}
and \texttt{t\_surj}.

For every valid \(P\), choose the point and
exponential of PUN22.  For invalid additive maps,
assign a fixed point of the same nonempty class
and the constant identity-valued map.
This defines the source's total functions on
\emph{all} additive maps.  Their axioms concern
valid periods only, so these extra values impose
no condition.  HUF14 gives both coverage and
disjointness for these choices.  To see coverage
literally, transport the period supplied by HUF6
through the integral isometry to the selected
representative; PUN23 then gives \(u_i(P)=y\).

For any open \(B\subseteq\mathbb C^m\) and any
holomorphic period family, the chosen pointwise
functions satisfy all the hypotheses of HUF12.
Thus they have the \emph{same} holomorphic maps
and joint exponential maps required by
\texttt{FamilyUnif.hol}; no separate choice for
each base or family is made.  Their derivative
clauses already belong to \(\mathcal U\), by HUF1.
The remaining pointwise exponential, action,
Weil and marking fields are likewise exactly
those of \(\mathcal U\).
\end{proof}
% END HUF15

% BEGIN HUF16
\paragraph{HUF16. The two exact remaining fiber and lift fields.}
\label{claim:HUF16}
For the single family chosen in HUF15, the
original \texttt{FamilyUnif.fibre} and
\texttt{FamilyUnif.lift} fields hold with their
printed half-lattice and full-\(A\) congruences.

\begin{proof}
PUN23 says that \(u_i(P)=u_i(P')\) is equivalent
to an \(R\)-linear isometry \(\gamma:L_i\to L_i\)
and a complex-linear \(h\) with
\(P'\gamma=hP\), and with
\(\gamma t_i a-t_i a\in ML_i\) on
\(\bar{\mathfrak b}L\) and in \(AL_i\) on
all of \(L\).
Surjectivity of \(t_i\) modulo \(N\) transfers
these statements to the exact source lattices.
For clarity, an element of
\(\bar{\mathfrak b}L_i\) is a finite sum
\(\sum_j b_jx_j\) with \(b_j\in\bar{\mathfrak b}\).
Write \(x_j=t_i a_j+N z_j\).
Then this element equals \(t_i(\sum_j b_ja_j)\)
modulo \(N\bar{\mathfrak b}L_i\).
Since \(M\mid N\), its displacement under
\(\gamma\) is zero modulo \(M\).
The analogous argument for all of \(L_i\)
uses \(A\mid N\).  Conversely \(t_i\) maps
each required source lattice into the indicated
target lattice, so the target congruences imply
the original congruences.  This proves both
directions of the literal \texttt{fibre} field.

For the lift field, take any holomorphic
\(y:B\to X_{o,\sigma}\) and \(b_0\in B\).
HUF15 gives \(i,P_0\) with \(u_i(P_0)=y(b_0)\),
together with its chosen exponential.
HUF8 continues these on an open
\(B'\ni b_0\), \(B'\subseteq B\).
At each \(b\), that continued uniformization
and the selected one at \(P(b)\) have identical
period and marked-lattice data.  PUN23 with
\(\gamma=h=1\) gives \(u_i(P(b))=y(b)\).
All the original open-set, base-point and
holomorphic-period requirements are retained.
\end{proof}
% END HUF16

\subsection{Simultaneous normalized theta frames}

% BEGIN HUF17
\paragraph{HUF17. Normalized periods force the original theta coordinates to be holomorphic.}
\label{claim:HUF17}
Assume the exact hypotheses of
\texttt{FamilyUnif.theta}, with an arbitrary supplied
integral symplectic basis, sign unit \(s\), and
matrices \(h(b),\tau(b)\).
On the nonempty open base they necessarily satisfy
\[
       s=-1,\qquad \tau=\tau^t,\qquad
       h(b)=A(b)^{-1},\qquad \tau(b)=A(b)^{-1}C(b),
 \tag{HUF.12}
\]
where \(A\) and \(C\) are the columns of the two
halves of the supplied period basis.
In particular both matrices are holomorphic,
although their holomorphy is not assumed in the
source theta field.

\begin{proof}
The source equations are \(hA=1\) and
\(hC=\tau\).  These are square \(g\)-by-\(g\)
matrices, so \(A,h\) are invertible and the
last two formulas follow.  Their entries are
holomorphic because the period columns are,
and matrix inversion is holomorphic where the
determinant is nonzero.

In the coordinates \(w=hz\), the positive
Hermitian form has a matrix \(S\).
Its imaginary part vanishes on the first half
of the integral basis, so \(S\) is real symmetric
and positive definite.  Set \(\tau=X+iY\).
The mixed pairing on that basis gives
\(-SY=s\,1\).  The stated positive definiteness
of \(Y\), and that of \(S\), imply \(s=-1\)
and \(S=Y^{-1}\); here \(g\geq1\), since the
original \(r\geq1\).  The pairing vanishes on
the second half too.  Its imaginary matrix is
\(\operatorname{Im}(\tau^tS\bar\tau)=X-X^t\),
so \(X=X^t\).  As \(Y=Y^t\), this proves
\(\tau=\tau^t\).
The sign is therefore the one appropriate to
our first-variable-linear Hermitian convention:
\(B'(a_i,b_j)=-\delta_{ij}\).
When the base is empty the desired theta
assertion has no conditions there, and needs
no sign conclusion.
\end{proof}
% END HUF17

% BEGIN HUF18
\paragraph{HUF18. One constant half-characteristic and its exact theta line.}
\label{claim:HUF18}
Under HUF17 there are constants
\(c_1,c_2\in\{0,\tfrac12\}^g\), depending on
the original marked lattice and the given
symplectic basis, such that the holomorphic
theta line on the period-torus family has
the same rigidified line as the actual
universal primitive line on each fiber.
Its section on the covering space is precisely
the source series
\(\operatorname{thetaChar}(\tau(b),c_1,c_2,h(b)z)\).

\begin{proof}
Write the symplectic basis as \(a_i,b_i\), and
use \(w=h(b)z\) as the covering coordinate.
The period lattice is
\(\mathbb Z^g+\tau(b)\mathbb Z^g\).
On the fixed integral lattice the primitive
identity, by the explicit Appell--Humbert
calculation in PUN20, forces
\[
          \alpha(x)=\exp(\pi iB'(x,w_0x))\in\{1,-1\}.
 \tag{HUF.13}
\]
This conclusion applies to the actual line:
its alternating form is \(B'\), since its
Weil formulas agree with (HUF.2) for every
positive integer.  The difference of two
integral forms with those values is divisible
by every such integer and hence zero.
The Hermitian form is then determined by
\(B'\) and the complex structure.

Choose \(c_{1i},c_{2i}\) in the stated two-element
set so that
\(\alpha(a_i)=e^{2\pi i c_{1i}}\) and
\(\alpha(b_i)=e^{-2\pi i c_{2i}}\).
The semicharacter law and the sign in HUF17
give, for \(m,n\in\mathbb Z^g\),
\[
 \alpha(m+\tau n)=
       (-1)^{m^tn}\exp(2\pi i(c_1^tm-c_2^tn)).
 \tag{HUF.14}
\]
All these signs are fixed by the lattice data;
none is chosen separately at different base points.

The holomorphic factors
\[
 j_\theta(m+\tau n,w)=
 \exp\!\left(2\pi i c_1^tm-\pi i n^t\tau n
                  -2\pi i n^t(w+c_2)\right)
 \tag{HUF.15}
\]
obey the lattice cocycle identity, by direct
multiplication and the integrality of the
cross terms.  They therefore define a
holomorphic line on the torus family, with
its canonical rigidification at \(w=0\).
On a fixed fiber put \(Y=\operatorname{Im}\tau\)
and \(H(w,v)=w^tY^{-1}\bar v\).
Multiplying the Appell--Humbert frame by the
quadratic gauge
\(\exp(-\pi w^tY^{-1}w/2)\) changes its factors
to (HUF.15).  Indeed the factor quotient is
\[
 \alpha(m+\tau n)\,
 \exp(-2\pi i n^tw-\pi i m^tn-\pi i n^t\tau n),
 \]
which is exactly (HUF.15) by (HUF.14).
This checks equality of the line itself and
its rigidification, not only of its Chern
class.  The gauge is used on a single fiber;
no holomorphy of \(Y\) in the base is asserted.

Reindexing the source theta series under
\(w\mapsto w+m+\tau n\) gives (HUF.15).
It is jointly holomorphic: on every compact
subset the smallest eigenvalue of \(Y\) is
bounded below by a positive constant, and
its summands and all differentiated summands
are bounded by a polynomial times
\(\exp(-c\|k+c_1\|^2+C\|k+c_1\|)\).
The resulting normally convergent series
therefore gives a holomorphic section of
this exact line.  It is not identically
zero on any fiber, as is also seen from
its nonzero Fourier coefficients.
\end{proof}
% END HUF18

% BEGIN HUF19
\paragraph{HUF19. Rigidified fiber isomorphisms glue over an arbitrary open base.}
\label{claim:HUF19}
Let \(\pi:T\to B\) be a proper holomorphic
submersion with connected compact fibers and
a zero section.  If two rigidified holomorphic
lines on \(T\) are isomorphic on every fiber,
then they have a unique global holomorphic
isomorphism respecting the rigidifications.
This applies to the theta line and actual
line of HUF18 on every original open
\(B\subseteq\mathbb C^m\).

\begin{proof}
Let \(Q\) be the Hom line between them.
On each fiber it is trivial, and its space
of holomorphic sections has dimension one,
because holomorphic functions on a compact
connected complex manifold are constant.
The line \(Q\) is flat over \(B\): locally
the submersion is a projection of complex
coordinate polydiscs, and \(Q\) is locally
free of rank one there.

Apply Grauert's constant-cohomology theorem
with degree zero.  It says that
\(E=\pi_*Q\) is a holomorphic line bundle
and that \(E\otimes\mathbb C(b)\simeq
H^0(T_b,Q_b)\) for every \(b\).
The evaluation map \(\pi^*E\to Q\) is
an isomorphism: on each fiber it evaluates
the constants in a trivial line and is
nonzero at every point.  Evaluating also
on the zero section gives an isomorphism
\(E\to e^*Q\).  The two rigidifications
identify \(e^*Q\) with \(\mathcal O_B\).
Their unit section therefore defines a
global nowhere-zero section of \(Q\),
normalized at zero.  It is the required
line isomorphism.  Another such isomorphism
has a ratio constant on each fiber and
equal to one at its zero point, so is
identical.

The precise named input is Grauert,
\emph{Ein Theorem der analytischen
Garbentheorie und die Modulr\"aume komplexer
Strukturen}, Publ.\ Math.\ IHES 5 (1960),
\S7, Satz 5, pp.\ 62--64, with the 1963
corrigendum in volume 16, pp.\ 35--36.
Its hypotheses here are a proper map,
a reduced base, a base-flat coherent sheaf
and constant \(h^0=1\), all just checked.
It gives local freeness, not a gratuitous
global trivialization: the latter is
provided by the rigidification argument
above.  No Stein, contractibility, or
simple-connectivity condition on \(B\)
is used.
\end{proof}
% END HUF19

% BEGIN HUF20
\paragraph{HUF20. The full simultaneous normalized theta-frame producer.}
\label{claim:HUF20}
For every open preconnected \(B\subseteq\mathbb C^m\),
every holomorphic period family, every original
symplectic basis and sign, and every \(h,\tau\)
in the printed theta hypotheses, suppose the
arbitrarily supplied \(y,e\) satisfy pointwise
\(\mathcal U\), \(\texttt{IsHolTo}(B,y)\), and
the printed joint holomorphy condition for \(e\).
Then the constants
of HUF18 satisfy the exact
\(\texttt{IsThetaFrame}\) predicate for the
actual universal rigidified line.
Thus the whole original
\texttt{exists\_thetaFrame\_family} conclusion
is supplied.

\begin{proof}
The given joint exponential map (supplied by HUF12
when applied to the chosen family) identifies the holomorphic
torus family with the analytic pullback of
the universal abelian scheme.  Indeed it
descends through the exact lattice, and
on each fiber is a bijective local
biholomorphism; its total differential is
invertible because it is the identity on
the base.  Thus the descended map is an
analytic isomorphism.  The torus family is
proper: its underlying real torus bundles
are locally products with a compact torus,
using a real basis of the holomorphic
period lattice.

HUF17 derives the holomorphic coordinates,
and HUF18 constructs the theta line and
its constants.  Pull back the actual
universal line to the torus family.
HUF19 gives the unique global rigidified
holomorphic isomorphism from the theta
line to it.  On the universal covering
\(B\times\mathbb C^g\), the theta line
has its canonical nonzero frame.
Transport this frame by that isomorphism,
and write its coefficient \(fr_j(b,z)\)
in each of the source's algebraic bundle
charts.

These coefficients are holomorphic on
the exact inverse images of the chart
opens, nowhere zero there, and obey
\[
 fr_j(b,z)=g_{jj'}(e_b(z))\,fr_{j'}(b,z).
 \tag{HUF.16}
\]
Because the isomorphism is rigidified,
their values at \(z=0\) equal the
printed rigidification section in
those same charts.  These verify
the first four \texttt{IsThetaFrame}
clauses with their actual transition
orientation.

The theta section of HUF18 transports
to a holomorphic section of the actual
line on each compact projective fiber.
Projective GAGA makes it an algebraic
section on that fiber.  In every
pullback chart its value is exactly
\(\operatorname{thetaChar}(\tau(b),c_1,c_2,h(b)z)
fr_j(b,z)\); this is the last clause,
including every point of the specified
fiber product.  That clause asks for
an algebraic section separately on
each fiber, not for algebraization
of the entire analytic family.
We have therefore proved exactly its
scope.  For \(B=\varnothing\), take
both characteristics to be zero and
arbitrary frame functions; all clauses
are vacuous.  We retain the source's
preconnectedness hypothesis throughout.
\end{proof}
% END HUF20

\subsection{Installation of the literal finite-family theorem}

% BEGIN HUF21
\paragraph{HUF21. The whole original \texttt{N1c} and \texttt{FamilyUnif}.}
\label{claim:HUF21}
For every original \(E,H,D,l,X_o,\sigma,g,\mathrm{ty}\)
with \(p+q\geq1\), \(\sigma(\rho)=\mathrm{rt}_0\),
and the given \texttt{IsTypeFn}, there exists a
\[
                       \texttt{FamilyUnif}\ X_o\ \sigma\ \mathrm{ty}.
 \tag{HUF.17}
\]
This is the literal \texttt{N1c} conclusion.
It retains the supplied \texttt{Moduli} record
and the named classical foundations stated above;
no individual-point or supplied-family hypothesis
remains in place of the family conclusion.

\begin{proof}
Take the finite positive family and the single
total maps \(u_i,e_{i,P}\) of HUF15.
Their genus records supply \texttt{sig},
\texttt{t\_inj} and \texttt{t\_surj}.
The complete pointwise \(\mathcal U\) supplies
\texttt{over\_\(\sigma\)}, \texttt{ex\_add},
\texttt{ex\_onto}, \texttt{ex\_ker},
\texttt{ex\_lin}, \texttt{ex\_weil},
\texttt{ex\_eta}, \texttt{ex\_etaA}
and both \texttt{ex\_diff} clauses.
HUF15 supplies \texttt{onto}, \texttt{disj}
and the universal \texttt{hol} field.
HUF16 supplies the exact \texttt{fibre}
and \texttt{lift} fields.
HUF20 supplies \texttt{theta}, for all
parameter dimensions, all open preconnected
bases, every prescribed symplectic basis
and every pair of normalizing matrices
allowed by the original statement.
These exhaust the source record.

Every choice of components and total maps
was made once, before the universal
period, base, matrix, basis and torsion
quantifiers.  All Weil assertions use
every positive integer with the original
negative sign.  The lattice maps are
integral modulo the exact \(N\);
the half-level and full-\(A\) conditions
have not been replaced by rational
or finite-index substitutes.
The case \(d=0\) was included in HUF9
and HUF11, and arbitrary \(m\), including
zero, is allowed throughout.
Thus (HUF.17) has precisely the original
quantifier order and scope.
\end{proof}
% END HUF21

\paragraph{Present boundary.}
HUF21 supplies the finite holomorphic family
that PUN11 previously had to assume.
It does not itself construct the representing
\texttt{Moduli} scheme.
The adjoining principal-chart theorem additionally
requires its own simultaneous normalized framed
points, the exact \texttt{IsUnf} dictionary and
the reduced-test-ring \texttt{FromTest} assertion.
Those are separate clauses, not consequences
of merely naming HUF21; their installation is
kept distinct from the family endpoint here.
